\section{Grounded hope and unresolved limits}
\label{sec:hope}

There are three demonstrated mechanisms that justify pursuing this program, and a complementary incentive hypothesis worth testing. First, universal reasoning about a small action boundary can replace an enormous family of claims about individual model intentions. Second, a shared prerequisite can sometimes block many possible harmful sequences; the finite graph and transaction examples demonstrate this mechanism exactly. Third, negative results can produce durable improvements: replay protection, current-state checks, independent receipts, and cumulative ledgers continue to work across many different proposal policies within their assumptions.

The persistent-allowance lemma identifies one way gains can accumulate: a checked constraint can remain in force while the system proposing actions changes. The specification-omission control identifies the corresponding danger: the process can also accumulate certificates for an inadequate contract. Progress requires retaining both the constraint and scrutiny of what it protects.

The caretaker hypothesis adds a further possibility: improve what systems propose while preserving controls on what they may execute. It remains an experiment to run, rather than an established fourth result.

These gains can support a constructive cycle. More capable systems help draft specifications, discover mistakes, and propose better implementations. Independent checks determine which improvements are retained. Useful work proceeds within the validated scope. Additional authority follows new evidence rather than the mere arrival of a more capable model. The cycle is an engineering program to test, not an already demonstrated route to arbitrary superintelligence.

\para{The largest open question is adequacy.} The world's transition system is not known. Software interfaces do not exhaust persuasion, social dependence, physical consequences, or actions by malicious human operators. Some valuable tasks may resist narrow specifications, and stronger systems may find omitted channels more effectively. We cannot currently assign a defensible universal $u_N$ to these uncertainties. The correct response is to retain the unknown and limit the claim.

\para{Control has costs and can be abused.} Concentrated authority can itself harm people. Decisions about legitimate scope require accountable institutions, contestability, privacy protections, and limits on the controller, as Section~\ref{sec:legitimacy} makes explicit. Different communities may disagree about which outcomes are acceptable. Our formalism cannot settle those disagreements. It can make explicit whose decisions a mechanism enforces and what evidence supports its technical promises.

\para{The program may not scale far enough.} If real systems cannot be usefully mediated, if capability gains invalidate attack coverage faster than it can be rebuilt, or if institutions cannot preserve practical revocation, this approach will be insufficient. Some failures will be visible in tests; others may remain undetected. We do not promise that every dangerous transition comes with an advance warning.

\para{Progress need not settle every dispute.} Better isolation, reliable execution budgets, validated recovery, and clearer evidence protect against ordinary errors and misuse as well as deliberate subversion. That makes the research worthwhile across a wide range of views about future AI. Skeptics help identify physical and institutional constraints; risk-focused researchers help pressure-test them. Neither needs to be dismissed for this work to proceed.
